\begin{lemma}[Removing a pinched weight from a Caccioppoli inequality]
Let $n\ge1$, let $s,K,c\in\R$ with $K\ge1$ and $c>0$.  Then there is a constant
$\gamma>0$, depending only on $n$, $s$, $K$ and $c$, with the following
property.

Let $x_0\in\R^n$, let $R,\Theta,\ell,\kappa_1,\kappa_2\in\R$ with $R>0$ and
$\kappa_1>0$ and $\kappa_2\le K\kappa_1$, write $B_\sigma=B_\sigma(x_0)$, let
$f\colon\R^n\to\R$, let $G\colon\R^n\to\R^n$ and let $w\colon\R^n\to\R$ be
measurable.  Assume
\begin{itemize}
\item[\textup{(H1)}] $f\in W^{1,2}_{\mathrm{loc}}(B_R)$ with weak gradient $G$
  (\noderef{MemW1pLocOn});
\item[\textup{(H2)}] $|f|\le\Theta$ almost everywhere on $B_R$;
\end{itemize}
Then for all radii $0<\rho_0<\rho_1<R$ the following two implications hold.
\begin{itemize}
\item[\textup{($+$)}] If, in addition,
  \begin{itemize}
  \item[\textup{(H3$^+$)}] $\kappa_1\le w\le\kappa_2$ almost everywhere on
    $B_R\cap\{f>\ell\}$, and
  \item[\textup{(H4$^+$)}] for every $\zeta\in C^\infty(\R^n)$ with compact
    support and $\operatorname{tsupport}\zeta\subseteq B_R$,
    \[
    \int_{B_R\cap\{f>\ell\}}w^{\,s}\|G\|^2\zeta^2\,dx
    \ \le\ c\int_{B_R}w^{\,s}\bigl((f-\ell)^+\bigr)^2|\nabla\zeta|^2\,dx ,
    \]
  \end{itemize}
  then
  $\mathrm{DeGiorgiLevelBound}\ n\ \gamma\ f\ G\ x_0\ \ell\ \rho_0\ \rho_1$.
\item[\textup{($-$)}] If, in addition,
  \begin{itemize}
  \item[\textup{(H3$^-$)}] $\kappa_1\le w\le\kappa_2$ almost everywhere on
    $B_R\cap\{f<\ell\}$, and
  \item[\textup{(H4$^-$)}] for every such $\zeta$,
    \[
    \int_{B_R\cap\{f<\ell\}}w^{\,s}\|G\|^2\zeta^2\,dx
    \ \le\ c\int_{B_R}w^{\,s}\bigl((\ell-f)^+\bigr)^2|\nabla\zeta|^2\,dx ,
    \]
  \end{itemize}
  then
  $\mathrm{DeGiorgiLevelBound}\ n\ \gamma\ (-f)\ (-G)\ x_0\ (-\ell)\ \rho_0\
  \rho_1$.
\end{itemize}
Both predicates are to be read through \noderef{DeGiorgiLevelBound}.  Here
$(t)^+=\max\{t,0\}$,
$|\nabla\zeta|^2=\sum_i(\partial_i\zeta)^2$, and the two displayed predicates
are the two instances of the paper's condition (2.1) of the thesis that the De Giorgi
lemmas of \S2 consume.

\medskip
\noindent\textbf{What this lemma is for, and why the quantifier order is this
one.}  The paper uses one computation three times in \S3: a Caccioppoli
inequality carrying the weight $w^{p/2-1}$, $w=\varepsilon+|\nabla u|^2$, is
converted into a condition of the form (2.1) of the thesis by cancelling the weight
against its own two-sided bound and then specialising the test function to the
standard cut-off.  At thesis l.~884--885 the paper says exactly this: ``If
$\ell\in[\mu/2,\mu]$ \emph{on the set of integration} in (3.4) of the thesis we have
that $2\ell\ge|D_ku|\ge\ell$.  Thus one can easily cancel the $w^{p/2-1}$ as we
did to obtain (3.9) of the thesis for the case $\varepsilon>\lambda^2$.''  The lemma
is that computation, abstracted in the function as \S2's lemmas are, so that it
can be stated once and instantiated three times.

The pinch is attached to the \emph{set of integration} and not to the ball,
because that is the only form two of the three instances have: at
$\varepsilon\ge\lambda^2$ the paper's $w$ is pinched on all of $B_R$, but for
levels $\ell\in[\mu/2,\mu]$ it is pinched only where $|D_ku|$ is of order $\mu$.
A bound valid on the whole ball restricts to the set of integration, so the
whole-ball instances are special cases and nothing is lost.

$K$ is bound \emph{before} the existential for $\gamma$ and $\kappa_1,\kappa_2$
\emph{after} it, which is forced: the weight is removed from the left of
(H4$^\pm$) at the cost $K^{|s|}\kappa_1^{-s}$ and reinstated on the right at the
cost $K^{|s|}\kappa_1^{\,s}$, and the two powers of $\kappa_1$ cancel exactly,
so the composite constant is $c\,K^{2|s|}$ times a factor depending only on the
cut-off.  It therefore carries no trace of $\kappa_1$ or $\kappa_2$ separately,
which is what lets one $\gamma$ serve instances with wildly different pinch
levels, and it is what the paper's ``$C(n,p,\gamma_1/\gamma_0)$'' asserts.

The two sign cases are stated as \emph{separate} implications, each carrying its
own pinch and its own Caccioppoli inequality on the set its own conclusion
integrates over.  They are genuinely different, because the second predicate
integrates over $\{-f>-\ell\}=\{f<\ell\}$, which is \emph{disjoint} from
$\{f>\ell\}$, so (H3$^+$) and (H4$^+$) say nothing about it.  Decoupling them is
not a stylistic choice but forced by the consumers: in the regime of the paper's
(3.23) of the thesis, where the weight is pinched only where $|D_ku|$ is of order
$\mu$, the $+$ pinch is available for levels $\ell\in[\mu/2,\mu]$ and the $-$
pinch only for $\ell\in[-\mu,-\mu/2]$ --- never both at the same $\ell$ --- so a
conjunctive form with both pinches as common hypotheses would be unusable there.
Both halves are nevertheless needed, because Lemma~2.5 of the paper asks for a
family of (2.1) of the thesis instances for $f$ at high levels and a family for $-f$
at low ones.  (That lemma is not cited here: the present statement is independent
of it, and naming it is commentary.)

The bound (H2) is hypothesised rather than derived.  In the paper's concrete
setting the weight is built from the very function being bounded, so the upper
pinch $w\le\kappa_2$ does bound $|f|$; once the lemma is abstracted in $f$ that
link is severed, and indeed $w\equiv1$ satisfies the pinch with
$\kappa_1=\kappa_2=K=1$ while $f$ is arbitrary.  (H2) is stated on $B_R$ rather
than on the inner ball because both clauses of
\noderef{DeGiorgiLevelBound} look at $B_{\rho_1}$ --- the first bounds
$(f-\ell)^+$ there and the second takes its essential supremum there --- and the
conclusion is quantified over all $\rho_1<R$.
\end{lemma}
\begin{proof}
\emph{Step 1: the cut-off, and the constant $\Lambda$.}  Let $\rho_{1/8}$ be the
standard mollifier of the Preamble in dimension one at scale $\tfrac18$ ---
smooth, nonnegative, supported in $[-\tfrac18,\tfrac18]$, of total integral $1$
--- and put $\psi_\flat=\mathbf1_{(-\infty,1/2]}*\rho_{1/8}$.  Then
$\psi_\flat\in C^\infty(\R)$ with
$\psi_\flat'=\mathbf1_{(-\infty,1/2]}*\rho_{1/8}'$ (Preamble: convolution with a
smooth compactly supported kernel, the indicator being locally integrable),
$0\le\psi_\flat\le1$ --- because $0\le\mathbf1_{(-\infty,1/2]}\le1$ and
$\rho_{1/8}\ge0$ with $\int\rho_{1/8}=1$, so the convolution is an average of
values in $[0,1]$ --- and $\psi_\flat=1$ on $(-\infty,\tfrac38]$ and
$\psi_\flat=0$ on $[\tfrac58,\infty)$.  For the latter two, write the
convolution in the form the Preamble item defines,
$(\mathbf1_{(-\infty,1/2]}*\rho_{1/8})(t)
=\int\mathbf1_{(-\infty,1/2]}(y)\rho_{1/8}(t-y)\,dy$, and substitute
$y=t-\sigma$, which is legitimate by the translation and reflection invariance
of Lebesgue measure (Preamble: invariance of Lebesgue measure under
translations and unimodular linear maps), to get
$\int\mathbf1_{(-\infty,1/2]}(t-\sigma)\rho_{1/8}(\sigma)\,d\sigma$.  Now for
$t\le\tfrac38$ every $\sigma$ with $|\sigma|\le\tfrac18$ has
$t-\sigma\le\tfrac12$, so the indicator is $1$ on the whole support of
$\rho_{1/8}$ and the integral is $\int\rho_{1/8}=1$; and for $t\ge\tfrac58$
every such $\sigma$ has $t-\sigma\ge\tfrac12$, so the indicator vanishes except
possibly at the single point $\sigma=t-\tfrac12$, a Lebesgue-null set, and the
integral is $0$.  Finally
$|\psi_\flat'|\le\|\rho_{1/8}'\|_{L^1(\R)}=:\Lambda<\infty$, an absolute
constant.  Given $0<\rho_0<\rho_1$ put
\[
\zeta_{\rho_0,\rho_1}(x)=\psi_\flat\Bigl(\frac{|x-x_0|-\rho_0}{\rho_1-\rho_0}\Bigr).
\]
Then $0\le\zeta\le1$; $\zeta=1$ on $\overline B_{\rho_0}$; $\zeta=0$ outside
$\overline B_{\rho_0+\frac58(\rho_1-\rho_0)}$, so $\operatorname{supp}\zeta$ is
compact and contained in $B_{\rho_1}$; $\zeta$ is smooth, because
$x\mapsto|x-x_0|$ is smooth away from $x_0$ and $\zeta$ is constant on the
neighbourhood $B_{\rho_0}$ of $x_0$; and
$|\nabla\zeta|\le\Lambda/(\rho_1-\rho_0)$, with $\nabla\zeta$ supported in
$B_{\rho_1}$.

\emph{Step 2: the pinch.}  If $0<\kappa_1\le w_1\le\kappa_2$,
$\kappa_1\le w_2\le\kappa_2$ and $\kappa_2\le K\kappa_1$ with $K\ge1$, then for
\emph{every} real $t$
\[
w_1^{\,t}\ \le\ K^{|t|}\,w_2^{\,t} .
\tag{2.1$'$}
\]
Indeed $w_1/w_2\le\kappa_2/\kappa_1\le K$ and $w_2/w_1\le K$; if $t\ge0$ then
$(w_1/w_2)^t\le K^t=K^{|t|}$ by monotonicity of $\sigma\mapsto\sigma^t$ on
$(0,\infty)$, and if $t<0$ then $(w_1/w_2)^t=(w_2/w_1)^{-t}\le K^{-t}=K^{|t|}$;
multiplying by $w_2^{\,t}>0$ gives (2.1$'$).  Two consequences are used, both
with $w_1$ or $w_2$ equal to the constant $\kappa_1$: for every $v$ with
$\kappa_1\le v\le\kappa_2$,
\[
v^{\,s}\ \le\ K^{|s|}\kappa_1^{\,s}
\qquad\text{and}\qquad
\kappa_1^{\,s}\ \le\ K^{|s|}v^{\,s} .
\tag{2.2$'$}
\]
The first is (2.1$'$) with $w_1=v$, $w_2=\kappa_1$, the second with
$w_1=\kappa_1$, $w_2=v$, both legitimate since $\kappa_1\le\kappa_1\le\kappa_2$.
Note that no case distinction on the sign of $s$ is needed, which is the point
of (2.1$'$): the paper's two exponent regimes $p/2-1>0$ and $p/2-1<0$ are
covered at once, by $|s|$.

\emph{Step 3: the constant.}  Put
\[
\gamma=c\,K^{2|s|}\bigl(1+\Lambda^2\bigr) .
\]
Then $\gamma>0$, since $c>0$, $K^{2|s|}>0$ and $1+\Lambda^2\ge1$; and $\gamma$
depends only on $n$, $s$, $K$ and $c$, because $\Lambda$ is an absolute
constant.  (The factor $1+\Lambda^2$ rather than $\Lambda^2$ only avoids having
to decide whether $\Lambda>0$; since $\Lambda^2\le1+\Lambda^2$ the estimate
below is unaffected.)

\emph{Step 4: the $+$ implication.}  Everything below is proved for an
arbitrary admissible sextuple, and it is worth saying exactly which, because
Step 5 reuses it under a substitution: \emph{let $\Theta,\ell,\kappa_1,\kappa_2$
be reals with $\kappa_1>0$ and $\kappa_2\le K\kappa_1$, let $w$ be measurable,
and let $(F,\Gamma)$ be a pair with $F\in W^{1,2}_{\mathrm{loc}}(B_R)$ of weak
gradient $\Gamma$ and $|F|\le\Theta$ a.e.\ on $B_R$, satisfying the pinch on
$B_R\cap\{F>\ell\}$ and the Caccioppoli inequality (H4$^+$) with $F,\Gamma$ in
place of $f,G$; then
$\mathrm{DeGiorgiLevelBound}\ n\ \gamma\ F\ \Gamma\ x_0\ \ell\ \rho_0\ \rho_1$
for all $0<\rho_0<\rho_1<R$.}  Nothing else about $F$, $\Gamma$ or $\ell$ is
used; in particular neither (H3$^-$) nor (H4$^-$) appears below.  We write the
argument with $(F,\Gamma)=(f,G)$, which is the $+$ implication of the statement.
Fix $0<\rho_0<\rho_1<R$, assume (H3$^+$) and (H4$^+$), and put
$\zeta=\zeta_{\rho_0,\rho_1}$ of Step 1; it is smooth, compactly supported, and
$\operatorname{tsupport}\zeta\subseteq\overline
B_{\rho_0+\frac58(\rho_1-\rho_0)}\subseteq B_{\rho_1}\subseteq B_R$, so it is
admissible in (H4$^\pm$).  Write
\[
E=B_R\cap\{f>\ell\},\qquad E_1=B_{\rho_1}\cap\{f>\ell\},\qquad
E_0=B_{\rho_0}\cap\{f>\ell\},\qquad \Theta_\ell=\Theta+|\ell| .
\]

\textup{(4a)} \emph{The first clause of \noderef{DeGiorgiLevelBound}.}  It asks
for a $C$ with $(f-\ell)^+\le C$ almost everywhere on $B_{\rho_1}$.  First,
$\Theta\ge0$.  This is not free --- $\Theta$ is an arbitrary real in the
statement --- but it follows from (H2), because $|B_R|=\omega_nR^n>0$ (Preamble:
the volume of a ball, with $R>0$), so the set on which (H2) asserts
$|f|\le\Theta$ has positive measure and is in particular nonempty; at any of its
points $0\le|f|\le\Theta$.  Hence $\Theta_\ell=\Theta+|\ell|\ge0$.  Now by (H2)
and $B_{\rho_1}\subseteq B_R$ we have $|f|\le\Theta$ a.e.\ on $B_{\rho_1}$, so
a.e.\ there $f-\ell\le|f|+|\ell|\le\Theta+|\ell|=\Theta_\ell$, and with
$0\le\Theta_\ell$ this gives $(f-\ell)^+\le\Theta_\ell$.  Take
$C=\Theta_\ell$.

\textup{(4b)} \emph{Integrability.}  The closed ball $\overline B_{\rho_1}$ is
compact and contained in $B_R$, because $\rho_1<R$; so by (H1)
(\noderef{MemW1pLocOn}, whose second clause asserts exactly
$\|G\|\in L^2$ of every compact subset of $B_R$) the function $\|G\|^2$ is
integrable on $\overline B_{\rho_1}$, hence on every measurable subset of it.
Also $|B_{\rho_1}|=\omega_n\rho_1^{\,n}<\infty$ (Preamble: the volume of a
ball).  Three integrands are therefore integrable on the sets indicated:
\begin{itemize}
\item $\|G\|^2\zeta^2$ on $E$, because $0\le\zeta\le1$ and $\zeta=0$ off
  $B_{\rho_1}$, so $\|G\|^2\zeta^2\le\|G\|^2\mathbf1_{B_{\rho_1}}$, which is
  integrable; this is the domination criterion of the Preamble
  (\texttt{MeasureTheory.Integrable.mono'}), the measurability of
  $\|G\|^2\zeta^2$ coming from that of $G$ --- itself part of (H1), since
  $L^2$ membership includes almost-everywhere strong measurability --- and the
  continuity of $\zeta$;
\item $w^{\,s}\|G\|^2\zeta^2$ on $E$, because by (H3$^+$) and the first half of
  (2.2$'$) we have $0\le w^{\,s}\le K^{|s|}\kappa_1^{\,s}$ a.e.\ on $E$, so this
  integrand is dominated a.e.\ on $E$ by
  $K^{|s|}\kappa_1^{\,s}\|G\|^2\mathbf1_{B_{\rho_1}}$; $w^{\,s}$ is measurable
  because $w$ is;
\item $w^{\,s}\bigl((f-\ell)^+\bigr)^2|\nabla\zeta|^2$ and
  $\bigl((f-\ell)^+\bigr)^2|\nabla\zeta|^2$ on $B_R$: both vanish off
  $E_1$, since $((f-\ell)^+)^2=0$ where $f\le\ell$ and $\nabla\zeta=0$ off
  $B_{\rho_1}$, and on $E_1$ they are bounded a.e.\ by
  $K^{|s|}\kappa_1^{\,s}\Theta_\ell^2\Lambda^2(\rho_1-\rho_0)^{-2}$ resp.\
  $\Theta_\ell^2\Lambda^2(\rho_1-\rho_0)^{-2}$, by (H3$^+$), (2.2$'$), (4a) and
  Step 1; a bounded measurable function supported in a set of finite measure is
  integrable.
\end{itemize}

\textup{(4c)} \emph{Removing the weight on the left.}  Almost everywhere on $E$,
the second half of (2.2$'$) applied with $v=w(x)$ --- legitimate by (H3$^+$) ---
gives $\kappa_1^{\,s}\le K^{|s|}w^{\,s}$, hence, multiplying by
$\kappa_1^{-s}\|G\|^2\zeta^2\ge0$,
\[
\|G\|^2\zeta^2\ \le\ K^{|s|}\kappa_1^{-s}\,w^{\,s}\|G\|^2\zeta^2 .
\]
Both sides are integrable on $E$ by (4b), so integrating the inequality over $E$
(Mathlib: \texttt{MeasureTheory.integral\_mono\_ae}, for two integrable
functions ordered almost everywhere) gives
\[
\int_E\|G\|^2\zeta^2\,dx\ \le\ K^{|s|}\kappa_1^{-s}\int_Ew^{\,s}\|G\|^2\zeta^2\,dx .
\]

\textup{(4d)} \emph{Removing the weight on the right.}  Almost everywhere on
$B_R$,
\[
w^{\,s}\bigl((f-\ell)^+\bigr)^2|\nabla\zeta|^2
\ \le\ K^{|s|}\kappa_1^{\,s}\bigl((f-\ell)^+\bigr)^2|\nabla\zeta|^2 :
\]
on $E$ this is the first half of (2.2$'$) with $v=w(x)$, multiplied by
$((f-\ell)^+)^2|\nabla\zeta|^2\ge0$; and off $\{f>\ell\}$ both sides are $0$,
because $((f-\ell)^+)^2=0$ there.  Integrating over $B_R$, both sides being
integrable by (4b),
\[
\int_{B_R}w^{\,s}\bigl((f-\ell)^+\bigr)^2|\nabla\zeta|^2\,dx
\ \le\ K^{|s|}\kappa_1^{\,s}\int_{B_R}\bigl((f-\ell)^+\bigr)^2|\nabla\zeta|^2\,dx .
\]

\textup{(4e)} \emph{The chain.}  Combining (4c), (H4$^+$) and (4d), and using
$\kappa_1^{-s}\kappa_1^{\,s}=1$ and $K^{|s|}K^{|s|}=K^{2|s|}$,
\[
\int_E\|G\|^2\zeta^2\,dx
\ \le\ K^{|s|}\kappa_1^{-s}\cdot c\cdot K^{|s|}\kappa_1^{\,s}
\int_{B_R}\bigl((f-\ell)^+\bigr)^2|\nabla\zeta|^2\,dx
\ =\ c\,K^{2|s|}\int_{B_R}\bigl((f-\ell)^+\bigr)^2|\nabla\zeta|^2\,dx .
\]
This is the exact cancellation the statement block announces: no power of
$\kappa_1$ survives.

\textup{(4f)} \emph{The left-hand side dominates the one the conclusion names.}
Since $\zeta=1$ on $\overline B_{\rho_0}\supseteq E_0$, the integrand
$\|G\|^2\zeta^2$ agrees with $\|G\|^2$ on $E_0$; it is nonnegative, integrable
on $E$ by (4b), and $E_0\subseteq E$.  Hence (Mathlib:
\texttt{MeasureTheory.setIntegral\_mono\_set}, for a nonnegative function
integrable on the larger set)
\[
\int_{E_0}\|G\|^2\,dx=\int_{E_0}\|G\|^2\zeta^2\,dx\ \le\ \int_E\|G\|^2\zeta^2\,dx .
\]

\textup{(4g)} \emph{The right-hand side is the one the conclusion names.}  The
integrand $((f-\ell)^+)^2|\nabla\zeta|^2$ vanishes off $E_1$ and satisfies
$|\nabla\zeta|^2\le\Lambda^2(\rho_1-\rho_0)^{-2}$, so
\[
\int_{B_R}\bigl((f-\ell)^+\bigr)^2|\nabla\zeta|^2\,dx
\ \le\ \frac{\Lambda^2}{(\rho_1-\rho_0)^2}\int_{E_1}\bigl((f-\ell)^+\bigr)^2\,dx .
\]
Write $\mathcal S=\operatorname{ess\,sup}_{B_{\rho_1}}\bigl((f-\ell)^+\bigr)^2$,
which is finite by (4a) --- it is at most $\Theta_\ell^2$ --- and nonnegative,
the integrand being nonnegative.  Almost everywhere on $E_1\subseteq B_{\rho_1}$
we have $((f-\ell)^+)^2\le\mathcal S$, so, integrating against the constant
$\mathcal S$ over the finite-measure set $E_1$ (Mathlib:
\texttt{MeasureTheory.setIntegral\_mono\_on} against a constant, and
\texttt{MeasureTheory.setIntegral\_const}),
\[
\int_{E_1}\bigl((f-\ell)^+\bigr)^2\,dx\ \le\ \mathcal S\,|E_1| .
\]

\textup{(4h)} \emph{Conclusion of the $+$ clause.}  Chaining (4f), (4e) and
(4g),
\[
\int_{B_{\rho_0}\cap\{f>\ell\}}\|G\|^2\,dx
\ \le\ \frac{c\,K^{2|s|}\Lambda^2}{(\rho_1-\rho_0)^2}\,\mathcal S\,
\bigl|B_{\rho_1}\cap\{f>\ell\}\bigr|
\ \le\ \frac{\gamma}{(\rho_1-\rho_0)^2}\,\mathcal S\,
\bigl|B_{\rho_1}\cap\{f>\ell\}\bigr| ,
\]
the last step because $c K^{2|s|}\Lambda^2\le cK^{2|s|}(1+\Lambda^2)=\gamma$ and
$\mathcal S|E_1|\ge0$.  Together with (4a) this is precisely the pair of clauses
of $\mathrm{DeGiorgiLevelBound}\ n\ \gamma\ f\ G\ x_0\ \ell\ \rho_0\ \rho_1$
(\noderef{DeGiorgiLevelBound}): its first clause is (4a) and its second is the
display just obtained, once $\mathcal S$ is identified with the factor the
definition names.  Those two are the same number but not the same expression:
(4g) defined $\mathcal S=\esssup_{B_{\rho_1}}\bigl((f-\ell)^+\bigr)^2$, the
essential supremum \emph{of the square}, while clause (2) names
$\bigl(\esssup_{B_{\rho_1}}(f-\ell)^+\bigr)^2$, the square \emph{of the}
essential supremum.  They agree because $(f-\ell)^+\ge0$ pointwise and
$t\mapsto t^2$ is nondecreasing on $[0,\infty)$, so the essential supremum
commutes with it, and because both are finite by (4a); the quantity being
nonnegative, no sign case arises.

\emph{Step 5: the $-$ implication.}  Fix $0<\rho_0<\rho_1<R$ and assume
(H3$^-$) and (H4$^-$).  Apply Step 4 with $f$ replaced by $-f$, $G$ by
$-G$, $\ell$ by $-\ell$, and $w$, $\kappa_1$, $\kappa_2$, $\Theta$ unchanged ---
that is, with $(F,\Gamma,\ell)=(-f,-G,-\ell)$ in the standalone implication Step
4 opens with.  Its four hypotheses on the substituted data are available:
\begin{itemize}
\item $-f\in W^{1,2}_{\mathrm{loc}}(B_R)$ with weak gradient $-G$, which is
  clause (1) of \noderef{WeakGradientLinear} applied with $n$, $q=2$, the open
  set $\Omega=B_R$ --- openness being that clause's only side condition, and the
  clause carrying no restriction on $|\Omega|$ --- the coefficients $a=-1$,
  $b=0$, and both pairs taken to be $(f_1,g_1)=(f_2,g_2)=(f,G)$, whose
  $W^{1,2}_{\mathrm{loc}}(B_R)$ membership is (H1).  Its conclusion is the
  membership of the pair
  $\bigl(x\mapsto(-1)f(x)+0\cdot f(x),\ x\mapsto(-1)G(x)+0\cdot G(x)\bigr)$,
  which is $(-f,-G)$ pointwise;
\item $|-f|\le\Theta$ a.e.\ on $B_R$, which is (H2);
\item the pinch on $B_R\cap\{-f>-\ell\}$: that set equals
  $B_R\cap\{f<\ell\}$, on which the pinch is (H3$^-$);
\item the Caccioppoli inequality for $(-f,-G)$ at the level $-\ell$: its
  left-hand integral is over $B_R\cap\{-f>-\ell\}=B_R\cap\{f<\ell\}$ with
  integrand $w^{\,s}\|-G\|^2\zeta^2=w^{\,s}\|G\|^2\zeta^2$, and its right-hand
  integrand carries $((-f-(-\ell))^+)^2=((\ell-f)^+)^2$; so the required
  inequality is exactly (H4$^-$), which is assumed in this implication.
\end{itemize}
Step 4 therefore yields
$\mathrm{DeGiorgiLevelBound}\ n\ \gamma\ (-f)\ (-G)\ x_0\ (-\ell)\ \rho_0\
\rho_1$, with the \emph{same} $\gamma$, since $\gamma$ was built from
$n,s,K,c$ alone and none of those changed.  Note that this is not a relabelling
of the $+$ implication: the two conclusions integrate over
$\{f>\ell\}$ and $\{f<\ell\}$, which are disjoint, so neither (H3$^+$),
(H4$^+$) nor the $+$ conclusion says anything about the other side --- which is
exactly why the statement keeps the two implications separate.  This completes
the proof.
\end{proof}
